% Gauss, Werke X/1 (1917), native127–128 upper, printed120–121.
\SourceLeaf{127}\NoSourceHead
\SourceGroup{IV}
\HistoricTitle{[ÜBER DAS REGELMÄSSIGE SIEBZEHNECK.]}
\NumberedItem{1}
\begin{center}\large\textsc{Pfaff} an \textsc{Gauss}.\end{center}
\begin{flushright}\{Helmstedt, 22. März 1802.\end{flushright}
\noindent\makebox[\linewidth][s]{. . . . . . . . . . . . . . . . . . . . . . . . . . .}\par
\noindent\hspace*{\parindent}\textsc{Langsdorfs} Compendium [*)] habe ich kürzlich bekommen,
aber noch nicht durchgesehen, es ist jezt beym Buchbinder. Prof. \textsc{Pfleiderer}, den man
wohl als den Repräsentanten der ächt geometrischen Methode in Deutschland ansehen kann,
ist, wie er in einem Brief an mich geäussert hat, nicht damit zufrieden. In eben diesem Brief
hat er mir auch eine geometrische Construction des 17-Ecks mitgetheilt, nach Ihren Formeln
oder nach denjenigen, die ich ihm als Auszug aus Ihrem ersten Msct. überschickt hatte.
Soviel ich mich erinnere, habe ich Ihnen das vorigemal diese Construction nicht communicirt,
sie ist mit seinen Worten folgende:

»Diametro $CH$ circuli dati ad extremum eius $C$ normalis ducatur $PN$ recta: a qua
abscindantur primum $CF=\tfrac14CH$; tum $FM,FN$ utraque $=FH$. Bifariam in punctis
$L,O$ secentur rectae $CM,CN$: et ad centrum $D$ circuli propositi ductis $LD,OD$ rectis,
ab perpendiculo $PN$ abscindantur $LP=LD,\ OQ=OD$. Ad rectam $CQ$ applicetur
rectangulum $(CSQ)$ aequale rectangulo sub $DC$ et $CP$, et excedens quadrato
(quod fit, si recta $PN$ ad $Q$ punctum normalis ducitur $QR=CP$, et super diametro
$DR$ circulus describitur secans rectam $CP$ in $S$). Denique centro $C$, intervallo $CS$,
descri-
\par\smallskip\noindent\rule{.25\linewidth}{.4pt}\par\smallskip
{\small\noindent\hspace*{\parindent}[*) K. Chr. v. \textsc{Langsdorf},
\textit{Anfangsgründe der reinen elementaren und höheren Mathematik u. s. w.,}
Erlangen 1802.]\par}
\clearpage
\SourceLeaf{128}\ArithmeticHead{121}{R}{ÜBER DAS REGELMÄSSIGE SIEBZEHNECK.}
\noindent batur circulus, secans propositum circulum in $B$. Erit $HB$ latus polygoni
regularis XVII laterum, circulo dato inscribendi«.
\begin{center}
\begin{tikzpicture}[x=.012cm,y=-.012cm,line width=.45pt,font=\small,
  aux/.style={dash pattern=on 1.3pt off 1.4pt,line width=.35pt}]
\draw (934,819) circle[radius=135];
\draw (798,375)--(798,1090);
\draw (798,819)--(1069,819);
\draw[aux] (360,752)--(798,752);
\draw[aux] (798,654)--(934,819);
\draw[aux] (798,922)--(934,819)--(983,763);
\draw[aux] (798,752)..controls(848,748)and(896,779)..(934,819);
\draw[aux] (934,819)..controls(938,797)and(934,768)..(932,747);
\draw[aux] (1045,719)..controls(1055,749)and(1063,780)..(1069,819);
\draw[aux] (414,752)..controls(418,730)and(427,701)..(435,683);
\draw[aux] (414,752)--(476,761);
\draw[aux] (649,784)--(701,792);
\draw (670,784)--(669,793);
\draw[aux] (873,811)--(934,819);
\draw[aux] (798,435)..controls(821,434)and(841,438)..(860,444);
\draw[aux] (757,541)..controls(789,551)and(816,564)..(838,583);
\draw[aux] (798,558)--(856,563);
\draw (798,482)--(802,482);
\draw (798,1020)--(803,1020);
\draw (798,558) circle[radius=3];
\fill (798,819) circle[radius=1.8];
\fill (1063,773) circle[radius=1.8];
\node[anchor=east,inner sep=1pt] at(798,432){$P$};
\node[anchor=east,inner sep=1pt] at(798,478){$M$};
\node[anchor=east,inner sep=1pt] at(798,554){$S$};
\node[anchor=east,inner sep=1pt] at(798,649){$L$};
\node[anchor=south east,inner sep=1pt] at(798,752){$F$};
\node[anchor=north east,inner sep=1pt] at(798,752){$Q$};
\node[anchor=east,inner sep=1pt] at(798,827){$C$};
\node[anchor=east,inner sep=1pt] at(798,925){$O$};
\node[anchor=east,inner sep=1pt] at(798,1020){$N$};
\node[anchor=north,inner sep=1pt] at(414,752){$R$};
\node[anchor=south west,inner sep=1pt] at(1063,773){$B$};
\node[anchor=west,inner sep=1pt] at(1069,815){$H$};
\node[anchor=north,inner sep=0pt] at(934,819){$D$};
\end{tikzpicture}
\end{center}
\noindent\hspace*{\parindent}Ich habe schon vor einiger Zeit eine Construction gefunden,
die ich Ihnen gelegentlich ein andermal mittheilen will.
.\,.\,.\,.\,.\,.\,.\,.\,.\,.\,.\,.

Behalten Sie ferner in freundschaftlichem Andenken
\begin{flushright}\begin{tabular}{c}Ihren ergebensten Freund\\J. Fr. \textsc{Pfaff}.\}\end{tabular}\end{flushright}
